\documentclass{article}
\usepackage{amsmath,amssymb,amsthm,hyperref}
\newtheorem{theorem}{Theorem}\newtheorem{proposition}[theorem]{Proposition}
\title{Planned causal cuts as independent commands with native marked semantics}
\date{16 November 2026}
\begin{document}\maketitle
\paragraph{Status, source, and purity.}
This candidate uses the reviewed \texttt{Causal\_Frontier\_Strategies.tex}:
variable-arity primitive polynomial grades, complete causal frontier tables,
exact local lumpability, and the scope-atomic stopping theorem are inputs.
It constructs an independent command double category and identifies the
precise native semantics of its action. The command action and finite-list
monoidal envelope are formal constructions; their native interpretation and
marker-descent boundary are checked separately. No original contentad
completion, general polynomial substitution compiler, or novelty claim is
made. Unlike the reserve history-interface construction, raw history trees
remain \emph{internal system data}. Native witness maps with fixed pure
endpoints therefore remain globular system maps, not new pure-interface maps.

\section{Independent command doctrine}
A pure one-scope profile is a triple $e=(A,U,B)$ of external inputs, recursive
states and external outputs. Tight maps are triples of ordinary functions
$(a,u,b):e\to e'$. There are only loose endocommands $k:e\rightsquigarrow e$,
with $k\in\mathbb N$. A square from $k$ at $e$ to $k'$ at $e'$ exists over
$(a,u,b)$ on \emph{both} sides precisely when $k=k'$; it is the specified
triple itself. Horizontal composition adds command counts, identities have
count zero, and vertical composition composes pure triples. The two square
compositions commute. Thus this is a strict double category without any
strategy, probability, witness or stopping-rule data in its interactions.
It is not asserted to contain all pure companions or to be an equipment.

For the symmetric monoidal doctrine take finite lists of pure profiles.
Tight maps are scope bijections with a pure triple in each matched component;
loose interactions are vectors of natural-number commands, one per scope.
Squares require the target vector to be the corresponding permutation of the
source vector and have the same tight map on their left and right sides.
Composition adds matched coordinates. Tensor concatenates lists and vectors;
its symmetry permutes scopes. This is a symmetric monoidal strict double
category $\mathbb{Cmd}$. A scope is indivisible by definition. No operation
merges two scopes or splits one correlated scope into independent components.

\section{Controlled native systems}
A one-scope system contains the entire primitive process
$f:A+U\to F_C(B+U)$, its full native context and content coefficients, and a
complete causal frontier strategy $\sigma$. At every raw nonterminal history
$h$, its normalized law $q_h$ has the prescribed primitive marginals. All
zero histories, unused content shapes, and actual coefficient arrows are
retained. Its global terminating grade is the full $R_UC$.

Additionally supply a complete cut plan $\mathcal P=(B_j)_{j\geq1}$. Each
$B_j$ is a bounded stopping-frontier rule on complete raw histories, always
stopping on earlier terminal histories. For every terminal raw history $t$,
its cut depths obey
\[
 0\leq d_1(t)\leq d_2(t)\leq\cdots\leq\operatorname{height}(t).
\]
The corresponding stopping rules must be pathwise nondecreasing on all raw
paths, including nonterminating and zero-probability paths, not merely on
realized terminal trees. Repeated equal depths and terminal absorption are
allowed; they add marks, not idle execution nodes. For each fixed $j$ a
uniform finite bound is specified; no common bound for all $j$ is required.
A finite supplied plan can be extended by repeating its last stopping rule;
the empty plan can be extended by immediate cuts. Choosing such an extension
is additional control data, not a recovered probabilistic fact.

A system also has a cursor $n\in\mathbb N$. Its evaluated native semantics
retains the entire global grade
\[
 \operatorname{Mark}_n(R_UC)
 =\{(t,e_1,\ldots,e_n):0\leq e_1\leq\cdots\leq e_n
                                \leq\operatorname{height}(t)\}.
 \tag{G}
\]
At every marked shape the content coefficient is the tensor of all original
node content factors of $t$, with its actual specified coherence maps.
The raw terminating histories, weights $P_\sigma(t)$ and node-product context
coefficients are those of the source theorem. The evaluated target is the
rowwise countable subprobability native context model of that source: maps
preserve measures by exact pushforward, context coefficients are products of
augmented node coefficients, and content arrows are the actual arity-preserving
arrows between full marked polynomial grades. Outputs are not asserted to be
normalized. Terminating
histories need not have finite support. This is not an implicit claim that
the finite-support writer model of the earlier wiring result is closed under
recursive execution. The body at $t$ uses exactly
$(t,d_1(t),\ldots,d_n(t))$. Unused marker tuples are not deleted.

A native system map has equal cursors, a source primitive comparison over the
pure triple, and its induced raw-history map $H$. It must satisfy
\[
 (r_h)_*q_h=q'_{Hh}\quad\text{for every }h,
 \qquad H\text{ preserves all planned stopping decisions }B_j.
 \tag{N}
\]
In particular the chosen stopping depths agree on every mapped history.
The actual original context/content maps are retained. The full expanded
grade arrow is the primitive-induced tree map with all $n$ markers retained.
No arbitrary extra map of expanded tree grades is admitted merely because it
preserves the evaluated body. Identities and composites satisfy (N), so these
objects and arrows form a carrier category over pure profiles. Finite lists
of systems, componentwise maps and scope permutations give the carrier for
$\mathbb{Cmd}$, with tensor list concatenation.

\begin{theorem}[Pure-interface-preserving command module]
Command $k$ acts by advancing the cursor from $n$ to $n+k$, retaining the
entire primitive process, strategy, plan and all $n+k$ global markers. On a
native map it retains its primitive arrows and extends their induced tree map
to all the marked shapes. Componentwise action defines a symmetric monoidal
loose right module over $\mathbb{Cmd}$. Identity commands add no marks;
a genuine immediate planned cut consumes command one and is not identity.
\end{theorem}
\begin{proof}
Condition (N) at every planned cut ensures that the same primitive map defines
a valid system map at every later cursor. Thus command action is a functor
on the pullback of the carrier and loose-arrow categories over pure profiles.
Advancing by $k$ and then $l$ gives exactly cursor $n+k+l$, the same complete
control data, and the same induced map at every global marked shape as
advancing by $k+l$. The unit is count zero. All action coherence is therefore
literal at the control-system interface; if parenthesized coefficient models
are used, their canonical nodewise monoidal constraints give the corresponding
pseudo presentation. Concatenation and permutation of separate scopes commute
with the action coordinatewise. The projection and action are strong symmetric
monoidal, giving the claimed module in the carrier/action formulation of
Libkind--Myers, Definition 4.1 and Explication 4.2.
\end{proof}
This is not the regular-hom action: the interactions are natural-number
commands rather than arbitrary native processes. It is nevertheless important
not to oversell this formal step. The substantive interpretation of a command
as cutting an actual correlated branching execution is the next result.

\begin{theorem}[Exact marked decomposition semantics]
Evaluating a planned system at cursor $n$ is natively isomorphic to its
successive prefix/suffix decompositions at the first $n$ planned barriers,
provided both sides retain the same full grade (G). These isomorphisms are natural
in the native maps (N), compatible with further segmentation and list tensor.
Compatibility with further commands means that both decompositions at the
new cursor agree under the source cut-associativity maps; it does not posit
a native transformation from one cursor evaluation to the next.
Forgetting marks yields a generally noninvertible content comparison instead.
\end{theorem}
\begin{proof}
Apply the source stopping-frontier theorem repeatedly. Every complete raw
terminal tree has a unique sequence of the specified prefixes and residual
forests. The decomposition includes the complete raw past and cannot split a
frontier draw. Multiplying the segment transition products gives exactly
$P_\sigma(t)$, without dividing by prefix probabilities. Prefix laws are
normalized, while the last suffix can be subprobabilistic. The node-product
context coefficient is the categorical product of the segment products,
canonically and naturally, because both sides have the same node projections.
Content factors are the same full addressed-node tensors and all ordered
markers are retained on both sides, including unused marks. Rebracketing cuts
does not erase a marker. Local lumpability and preservation of every planned
barrier give the source native transport theorem at each stage. Its products
and induced full-grade maps have the same primitive components on both
routes. Separate scope lists are treated componentwise, not as one joint draw.
\end{proof}

\paragraph{A controlled source section.}
The canonical plan stops at round $j$, or terminal earlier, for each $j$.
Every source native strategy comparison preserves depth, active/terminal
status and all raw history prefixes, so preserves this plan. Equipping a
source strategy with this plan and cursor zero is therefore fully faithful
on the declared primitive-induced strategy category, with unchanged pure
profiles and native maps. Arbitrary other plans are additional controls, and
an arbitrary native map need not preserve them. This is not full faithfulness
into the category of every possible native map between evaluated tree
processes: that larger category admits grade/raw maps not induced by a source
primitive comparison. The independent-strategy locus still carries all
original rowwise arity-preserving primitive comparisons by the source's
proved product-marginal calculation.

\section{Exact marker descent at a fixed raw execution interface}
Fix the same complete raw terminal set $T$, the same weights and actual node
coefficients, and the same unmarked shape map $\tau:T\to R_UC$. Let two plans
or cursor presentations have $m$ and $n$ markers, with body functions
\[
 b:T\to\operatorname{Mark}_m(R_UC),\qquad
 b':T\to\operatorname{Mark}_n(R_UC).
\]
For each unmarked tree $t$ of height $h$, let $M_m(t)$ be the finite set of
weakly increasing $m$-tuples in $\{0,\ldots,h\}$; its cardinality is
$\binom{h+m}{m}$. Consider native comparisons whose raw map and node
coefficient arrows are identities and whose global grade map lies over the
identity of the unmarked tree grade. This is the declared descent question for \emph{evaluated terminating native
processes}, not a classification of arbitrary primitive simulations or
plan-compatible strategy maps. The latter must still respect all raw prefix
histories, including divergent histories, zero histories and all future
planned cuts; terminal-body factorization does not certify that stronger
causal-control descent.

\begin{theorem}[Marker factorization and invertibility]
Such a native comparison exists if and only if $b'$ is constant on every
fibre of $b$ on \emph{all} raw terminal histories, including zero histories.
An invertible such comparison exists if and only if, for each unmarked $t$, the partial map on
realized marker tuples determined by that factorization extends to a
bijection $M_m(t)\to M_n(t)$. For nonempty tree grades this forces $m=n$;
then the extension criterion is exactly injectivity of the well-defined
partial map on realized tuples.
\end{theorem}
\begin{proof}
The native body equation for identity raw map is $\eta b=b'$. Thus the
factorization condition is necessary. If it holds, it prescribes a function
on every realized tuple. Extend it on unused tuples by any fixed target tuple,
for example the all-zero tuple, separately over each $t$. Such a tuple exists
even when $n=0$. At every global shape use the identity of its underlying
node content coefficient; source and target have that same coefficient of
$t$. This gives an arity-preserving full-grade arrow and the native body
square. Raw effects and context arrows are unchanged, proving sufficiency.

A native isomorphism with the declared identity raw map requires a full grade
bijection over each $t$; conversely such bijections with identity content
coefficient maps provide the inverse. Their restrictions to realized tuples
must be the prescribed partial map. Finite sets of the same cardinality admit
a bijective extension of a partial function exactly when that function is
injective. For any nonempty tree, $h\geq1$, and
$\binom{h+m}{m}$ is strictly increasing with $m$. Therefore equality of these
cardinalities forces $m=n$. The empty full tree grade is a vacuous exception.
\end{proof}
This shows why marker erasure is always a directed comparison but not an
isomorphism over the original tree grade when genuine marks remain. Inserting
one mark into an unmarked native process is possible at this interface exactly
when the planned cut depth depends only on its unmarked content tree, not on
finer raw witness distinctions.

\paragraph{Raw-dependent cut counterexample.}
Take one unary shape $c$, one loop state and one external output. At each root
or recursive row there are three raw witnesses: $r_0,r_1$ return the loop state
and $e$ returns the external output, all with shape $c$, with probabilities
$1/4,1/4,1/2$. Stop after round one when the first witness is $r_0$, otherwise
stop by round two, always stopping on terminal histories earlier. This is a
bounded causal whole-frontier rule. The two raw trees $(r_0,e)$ and $(r_1,e)$
have the same unmarked two-node content chain, but cut depths one and two.
Therefore no body-preserving grade insertion from $R_UC$ to its full cut grade
exists with identity raw trees. The action remains valid: it changes the
system's marked semantic grade, and never claims to be a natural
transformation from unmarked evaluation to marked evaluation.

\paragraph{Zero-time unit counterexample.}
For a stronger unit test, take $A$, $U$ and $B$ each singleton and use
two raw witnesses at every row, each of
probability $1/2$, one returning the loop and one returning the external
output, both of the same unary shape. There is exactly one terminating raw
witness of every positive length $n$, of weight $2^{-n}$. These weights are
pairwise distinct. Any native isomorphism must biject raw witnesses and
preserve their individual weights, hence fixes every length, even if its raw
map was not assumed to be identity. If stop at zero were the identity
interaction, the unit body square would therefore send every realized
$(t_n,0)$ to $t_n$. These already exhaust the full unmarked grade, leaving no
injective images for the additional global marked shapes.
Equivalently the preceding cardinality test fails over each unmarked tree.
Thus zero-time cutting and doing nothing are genuinely different operations
at this retained-content interface.

\section{Limits of the result}
The command module keeps original pure profiles rather than exposing raw
history as pure data, retains complete plans and strategy tables rather than
an almost-sure quotient, and provides exact marked recursive evaluation.
It does not admit arbitrary scope splitting, scope merging, scheduler change,
mutable shared state, or every native arrow of an expanded output grade.
Its tensor is a list of separate scope profiles, not a correlated synchronized
parallel operation. The controlled source section does not infer any
unspecified plan or conditional law from a marginal.

In particular this does not compile arbitrary original polynomial dependent
substitution. Such a compiler must specify a typed call interface, shared
full global grades across body fibres, actual aggregation under native shape
merges, and a transformation of complete strategies preserving causal scopes.
The existing routing and fair-bit scope-splitting counterexamples remain
load-bearing tests. An independently specified polynomial term/call doctrine
might supply that extension, but merely replacing commands with arbitrary
processes would revert to the regular-hom repackaging and settle nothing.
\end{document}
